\documentclass[11pt]{article}

\usepackage[a4paper,margin=2.5cm]{geometry}
\usepackage{amsmath,amssymb,amsthm}
\usepackage{physics}
\usepackage{graphicx}
\usepackage{caption}
\usepackage{subcaption}
\usepackage{booktabs}
\usepackage{siunitx}
\usepackage{xcolor}
\usepackage{listings}
\usepackage{hyperref}
\usepackage{float}
\usepackage{enumitem}
\usepackage{authblk}

\sisetup{separate-uncertainty=true}

\hypersetup{
  colorlinks=true,
  linkcolor=blue!50!black,
  citecolor=blue!50!black,
  urlcolor=blue!50!black
}

\lstdefinestyle{pystyle}{
  language=Python,
  basicstyle=\ttfamily\scriptsize,
  keywordstyle=\color{blue!60!black}\bfseries,
  commentstyle=\color{green!40!black}\itshape,
  stringstyle=\color{orange!70!black},
  numbers=left,
  numberstyle=\tiny\color{gray},
  stepnumber=1,
  numbersep=6pt,
  frame=single,
  breaklines=true,
  breakatwhitespace=false,
  showstringspaces=false,
  tabsize=4,
  columns=flexible,
  literate=
    {¹}{{\textsuperscript{1}}}1
    {×}{{$\times$}}1
    {φ}{{$\varphi$}}1
    {—}{{---}}1
    {⁴}{{\textsuperscript{4}}}1
    {₀}{{\textsubscript{0}}}1
    {₁}{{\textsubscript{1}}}1
    {₂}{{\textsubscript{2}}}1
    {ₐ}{{\textsubscript{a}}}1
    {ₓ}{{\textsubscript{x}}}1
    {ₘ}{{\textsubscript{m}}}1
}
\lstset{style=pystyle}

\newcommand{\eV}{\text{eV}}

\title{\Large Computational Models in Quantum and Statistical Physics: \\
Ten Numerical Studies from Random Walks to Hydrogenic Orbitals}

\author{Prepared for the BPhO Computational Physics Challenge 2026}
\date{August 2026}

\begin{document}

\maketitle

\begin{abstract}
This paper documents and analyses ten computational physics programs written for the BPhO Computational Physics Challenge 2026 (Quantum Mechanics strand), together with two extension programs. Each program is treated as an independent numerical experiment: a physical model is stated, the governing equations are derived from first principles, the numerical implementation is described, and the output generated by the program is presented and checked, wherever possible, against an analytical prediction or a known physical limit. The programs span a deliberately wide range of physics, from the classical statistics of a two-dimensional random walk and a hard-sphere Brownian-motion simulation, through the semiclassical Planck and Einstein models of thermal radiation and specific heat, to fully quantum-mechanical treatments of the photoelectric effect, the Bohr hydrogen atom, electron diffraction, the infinite square well and the uncertainty principle, photon-polarisation measurement, relativistic Compton scattering, and the three-dimensional hydrogenic wavefunctions. Every model is validated against at least one independent check -- a full-curve theoretical comparison, recovery of a known input parameter, a required limiting behaviour, or an independent numerical cross-check of a closed-form result -- and each is found to agree with the underlying physics to the precision expected of a double-precision implementation. The paper is intended to demonstrate the complete chain of reasoning from physical law to numerical result: theory, discretisation, implementation and interpretation are treated as inseparable parts of the same computational investigation.
\end{abstract}

\tableofcontents

% =====================================================================
\section{Introduction}
% =====================================================================

Quantum mechanics was built historically on a sequence of models that were each, in isolation, only partially correct: the equipartition-violating Planck spectrum, the semiclassical Bohr atom, the corpuscular photoelectric effect, and the wave-mechanical infinite square well. Each model captured a piece of the eventual quantum-mechanical picture while remaining formally distinct from the others -- some are frankly classical statistical arguments (the random walk), some are ``old quantum theory'' (Bohr, Einstein solids), and some are genuinely quantum-mechanical (the particle in a box, hydrogenic orbitals, Compton scattering treated relativistically). This collection of ten short programs, together with two interactive/verification extensions, reconstructs that sequence numerically. The purpose of each program is not merely to reproduce a textbook plot, but to start from the underlying physical assumptions, derive the specific equations that are then implemented, and use the resulting numbers or figures to test the model against a known result: the diffusive scaling law of a random walk, the Dulong--Petit limit, the correct ground-state energy of hydrogen, the Heisenberg bound, or the low- and high-voltage limits of an electron-diffraction pattern.

This paper is organised task-by-task, broadly following the order of increasing physical sophistication used in the challenge: statistical/classical models (Tasks~1--2), semiclassical radiation and solid-state models (Task~3), the historical quantum-mechanical results of Einstein, Bohr and de Broglie (Tasks~4--6), genuinely wave-mechanical quantum mechanics (Task~7 and its extension), quantum measurement in the context of polarisation and cryptography (Task~8), special-relativistic photon scattering (Task~9), and the full three-dimensional hydrogen wavefunction (Task~10). Section~\ref{sec:methods} describes the numerical infrastructure common to all ten programs. Sections~\ref{sec:task1}--\ref{sec:task10} treat each task in turn, and Section~\ref{sec:extensions} discusses the two extension programs (an interactive photoelectric-effect simulator and a numerical/analytic cross-check of the Heisenberg uncertainty principle) in more depth than their brief mentions in Sections~\ref{sec:task4} and \ref{sec:task7}. Section~\ref{sec:discussion} draws together recurring themes -- validation strategy, numerical method, and the idealisations common to all ten models -- before the conclusions in Section~\ref{sec:conclusions}. Full source code for every program is reproduced in the appendices.

Throughout, physical constants are taken from the 2018 CODATA recommended values as tabulated by NIST \cite{codata2018}, and are hard-coded into each program to double precision; no external constants library is used. All figures reproduced in this paper were generated by running the submitted programs exactly as supplied, using \texttt{Python} 3 with \texttt{NumPy} \cite{numpy} and \texttt{Matplotlib} \cite{matplotlib}; no data in this paper have been invented or adjusted, and any numerical values quoted in the text are taken directly from the programs' console output or read from the generated figures.

% =====================================================================
\section{Computational Methods}
\label{sec:methods}
% =====================================================================

All ten core programs share a common structure. Each begins by fixing the \texttt{Matplotlib} backend to the non-interactive \texttt{Agg} renderer and redirecting its configuration directory to a location next to the script, so that every program can be executed in a headless environment without a display and without requiring write access outside its own folder. Physical constants (the Planck constant $h$, the elementary charge $e$, the electron mass $m_e$, the speed of light $c$, the vacuum permittivity $\varepsilon_0$, the Boltzmann constant $k_B$ and the molar gas constant $R$) are defined once at SI precision near the top of each file and are not imported from any third-party constants package.

Numerical work is vectorised using \texttt{NumPy} array operations wherever the underlying mathematics permits it -- for example, the $n$-walker random walk of Task~1 and the Bragg-ring calculation of Task~6 evaluate all trajectories, or all diffraction orders, in a single array expression rather than a Python-level loop. Two programs (Task~2 and Task~10) additionally perform explicit control-flow -- collision detection in Task~2, and recurrence relations for the Laguerre and Legendre polynomials in Task~10 -- because these calculations are not naturally expressible as single vectorised operations. Figures are saved directly to Portable Network Graphics (PNG) files at $150$--$180$ dpi; Task~2 additionally produces an animated GIF using \texttt{Matplotlib}'s \texttt{FuncAnimation} and a Pillow writer. Every program prints at least one quantitative check (a root-mean-square displacement, a threshold frequency, a fitted lattice spacing, a comparison with the Heisenberg limit) to standard output; these printed values, reproduced verbatim from the actual program runs used to generate this paper's figures, are quoted throughout Sections~\ref{sec:task1}--\ref{sec:task10}.

No program in this collection solves a partial differential equation by finite differences or otherwise integrates the time-dependent Schr\"odinger equation numerically; the quantum-mechanical tasks (7, 9, 10) instead evaluate closed-form analytic solutions (the particle-in-a-box eigenstates, the relativistic Compton kinematics, and the hydrogenic wavefunctions) at arrays of sample points. This distinction is made explicit here because it determines what ``validation'' means for each task: where a program evaluates an exact analytic solution, the corresponding check in this paper is against an independent, differently-derived analytic result (or, for Task~7's extension, an independent numerical integration of the same analytic wavefunction), rather than against a converged numerical solution of a differential equation.

% =====================================================================
\section{Task 1 --- Two-Dimensional Random Walk}
\label{sec:task1}
% =====================================================================

\subsection{Physical context}
A particle undergoing repeated, uncorrelated displacements of fixed magnitude $s$ in random directions is the simplest discrete model of diffusive transport, and underlies the classical explanation of Brownian motion that Einstein and Smoluchowski developed from 1905 onwards \cite{einstein1905}. Task~1 implements this model directly: an ensemble of independent two-dimensional random walkers, each taking $N$ steps of fixed length $s$ in an angle drawn from a uniform distribution, and tests the resulting mean-squared displacement against the standard diffusive scaling law.

\subsection{Theory and derivation}
Let the $k$-th step of a walker be the vector $\vec{\delta}_k = s(\cos\theta_k, \sin\theta_k)$, with $\theta_k \sim U(0, 2\pi)$ drawn independently for each step and each walker. The position after $N$ steps is the partial sum
\begin{equation}
\vec{r}_N = \sum_{k=1}^{N} \vec{\delta}_k .
\end{equation}
The squared displacement is
\begin{equation}
r_N^2 = \vec{r}_N \cdot \vec{r}_N = \sum_{k=1}^N |\vec\delta_k|^2 + \sum_{k \ne l} \vec\delta_k \cdot \vec\delta_l
= N s^2 + s^2\sum_{k\ne l}\cos(\theta_k - \theta_l).
\end{equation}
Because the $\theta_k$ are independent and uniformly distributed on $[0,2\pi)$, $\langle \cos(\theta_k-\theta_l)\rangle = 0$ for $k \ne l$ (the expectation factorises over the two independent angles and $\langle\cos\theta\rangle=\langle\sin\theta\rangle=0$ for a uniform angle). Taking the ensemble average therefore eliminates every cross term, leaving the exact result
\begin{equation}
\label{eq:msd}
\langle r_N^2 \rangle = N s^2 .
\end{equation}
Equation~\eqref{eq:msd} is the discrete analogue of the diffusion law $\langle r^2(t)\rangle = 4Dt$ in two dimensions, with the number of steps $N$ playing the role of elapsed time. It holds for any step-length distribution with finite variance and any isotropic, independent, identically distributed set of step directions; the fixed step length used here is a special (and, for verification purposes, convenient) case, since then $\langle |\delta_k|^2\rangle = s^2$ exactly rather than only in expectation.

\subsection{Computational method}
The function \texttt{random\_walk} in \texttt{task1.py} draws all $N \times n_\text{walkers}$ direction angles from \texttt{numpy.random.default\_rng} in a single call, converts them to Cartesian step vectors, and forms the trajectory with \texttt{numpy.cumsum} along the step axis; an explicit origin row is concatenated onto the front of the resulting array so that \texttt{positions[:, k, :]} is the position after exactly $k$ steps. This is a fully vectorised implementation: no Python-level loop over steps or walkers is used. Two ensembles are generated with independent seeds: fifty walkers of $N=1000$ steps (with unit step length $s=1$) for the trajectory plot, and three thousand walkers of the same length for the statistical test of Eq.~\eqref{eq:msd}. The mean-squared displacement $\langle r^2\rangle$ is computed at every intermediate step number by averaging $\sum_i x_i^2+y_i^2$ over the walker axis, giving a full curve of $\langle r_N^2\rangle$ against $N$ rather than a single end-point value.

\subsection{Results and validation}
Figure~\ref{fig:task1} shows the fifty individual trajectories (left) and the ensemble-averaged mean-squared displacement compared with the theoretical line $Ns^2$ (right). The simulated curve tracks the theoretical prediction closely over the full range $N=0$ to $1000$, with the small residual fluctuations expected from a finite ensemble of $3000$ walkers.

\begin{figure}[H]
\centering
\includegraphics[width=\linewidth]{figures/task1.png}
\caption{Left: fifty independent two-dimensional random walks of $N=1000$ unit steps, all starting from the origin (black dot). Right: the ensemble-averaged mean-squared displacement $\langle r^2\rangle$ from $3000$ independent walkers (blue), compared with the theoretical prediction $Ns^2$ (black, dashed) of Eq.~\eqref{eq:msd}.}
\label{fig:task1}
\end{figure}

The program also prints a scalar end-point check: the root-mean-square net displacement after $N=1000$ steps of the $3000$-walker ensemble, compared with the theoretical value $s\sqrt{N}$:
\begin{quote}\ttfamily
RMS displacement (simulated)\phantom{x}: 31.434\\
RMS displacement (theory s*sqrt(N)): 31.623
\end{quote}
The theoretical value is $\sqrt{1000}=31.623$; the simulated value differs from it by $0.6\%$, well within the expected sampling fluctuation of $O(1/\sqrt{n_\text{walkers}}) \approx 1.8\%$ for a $3000$-walker ensemble. This close agreement, over the full curve and not just the final point, is the key validation of the random-walk model.

\subsection{Limitations}
The model assumes a fixed step length; in a physical Brownian system the mean free path and thus the effective step length would itself fluctuate. Correlations between successive step directions (present, for example, in Task~2's collision-driven motion) are entirely absent by construction, since each $\theta_k$ is drawn independently. The model is therefore best understood as establishing the purely statistical baseline -- diffusive $\langle r^2\rangle \propto N$ scaling from uncorrelated isotropic steps -- against which the more physically detailed hard-sphere simulation of Task~2 can later be compared.

% =====================================================================
\section{Task 2 --- Brownian Motion from Hard-Sphere Collisions}
\label{sec:task2}
% =====================================================================

\subsection{Physical context}
Task~2 replaces the abstract random walk of Task~1 with an explicit microscopic mechanism: a heavy particle of mass $M$ and radius $R$, initially at rest, is bombarded by $N$ light particles of mass $m$ and radius $r$ confined to a square box, and is deflected by momentum-conserving elastic collisions. This is closer to the physical picture Einstein and Perrin actually had in mind for Brownian motion -- a large colloidal particle jostled by the much more numerous molecules of the surrounding fluid -- than the purely statistical model of Task~1, at the cost of requiring an explicit collision-detection and impulse-resolution algorithm rather than a single closed-form scaling law.

\subsection{Theory: elastic collisions in two dimensions}
Consider two particles of mass $m_1, m_2$ with velocities $\vec v_1,\vec v_2$ colliding elastically, with $\hat n$ the unit vector along the line joining their centres at the moment of contact. Because the collision is frictionless, only the velocity components \emph{along} $\hat n$ can be changed by the impulse; the tangential components $\vec v_{1t} = \vec v_1 - (\vec v_1\cdot\hat n)\hat n$ and $\vec v_{2t}$ are unaffected. Writing $v_{1n} = \vec v_1\cdot \hat n$ and $v_{2n}=\vec v_2\cdot\hat n$ for the normal components, conservation of momentum and kinetic energy along $\hat n$ reduce to the standard one-dimensional elastic-collision result,
\begin{equation}
v_{1n}' = \frac{(m_1-m_2)v_{1n} + 2m_2 v_{2n}}{m_1+m_2}, \qquad
v_{2n}' = \frac{(m_2-m_1)v_{2n} + 2m_1 v_{1n}}{m_1+m_2},
\end{equation}
and the post-collision velocities are reassembled as $\vec v_1' = v_{1n}'\hat n + \vec v_{1t}$, $\vec v_2' = v_{2n}'\hat n + \vec v_{2t}$. This is exactly the calculation performed by \texttt{elastic\_collision\_normal} in \texttt{task2.py}. In the physically relevant limit $M \gg m$ used in the simulation ($M=200$, $m=1$), the big particle's normal-velocity update reduces to $v_{Mn}' \approx v_{Mn} + \tfrac{2m}{M}(v_{mn}-v_{Mn})$: each individual collision imparts only a small kick $O(m/M)$, consistent with the physical picture of Brownian motion as the accumulated effect of many small, individually negligible impacts.

\subsection{Computational method}
The simulation is a simple explicit (Euler) time-stepping scheme: at each timestep of duration $\dd t$, every particle is first advanted by $\vec v\,\dd t$; particles are then reflected specularly off the square walls of side $L$ if their coordinate has moved past $r$ (or $R$, for the big particle) from a wall; finally, collisions between the big particle and any small particle whose centre-to-centre separation has fallen below $R+r$ are resolved with the elastic-collision formula above, followed by an explicit de-overlap correction that pushes the small particle back onto the contact circle to prevent numerical interpenetration. Small-particle-small-particle collisions are not modelled -- the $N$ light particles pass through one another -- which is a deliberate simplification appropriate to a dilute-gas ``ideal bath'' picture, discussed further below.

Because a genuine kinetic-theory gas would re-randomise a molecule's direction only after travelling, on average, one mean free path, the small particles' headings are not randomised every frame. Instead the code defines a re-randomisation interval $\Delta t_\text{turn} = K_n\, r / s$, where $K_n$ is a dimensionless "Knudsen-like" input parameter and $s$ the (constant) small-particle speed; this converts the stated mean free path $K_n r$ into a time interval using the constant particle speed. Every small particle's velocity direction is redrawn uniformly at random whenever this interval elapses, while the speed $s$ is preserved. This is the only place in the ten programs where a phenomenological, rather than first-principles, timescale is introduced: the exact distribution of free-path lengths in a real dilute gas is exponential (Poissonian), whereas this implementation uses a single deterministic re-randomisation period for all $N$ particles simultaneously.

\subsection{Results}
Figure~\ref{fig:task2} shows three frames extracted from the $200$-frame animation \texttt{task2\_brownian\_motion.gif} produced by the program (the animation itself cannot be reproduced on a printed page). The big particle (red disk, radius $R=4$, versus small-particle radius $r=0.6$ and mass ratio $M/m=200$) is displaced from its initial position at the centre of the box under the cumulative effect of collisions with the surrounding bath of $150$ small particles, and by the final frame has traced out a visibly irregular path (red trail), qualitatively the erratic, direction-changing trajectory characteristic of Brownian motion.

\begin{figure}[H]
\centering
\includegraphics[width=\linewidth]{figures/task2_brownian_frames.png}
\caption{Three frames (1, 101 and 200 of 200) extracted from the animated hard-sphere Brownian-motion simulation of Task~2: $N=150$ small particles (blue dots, $r=0.6$) in a box of side $L=100$, and one heavy particle (red disk, $R=4$, $M/m=200$), initially at rest at the box centre. The red curve in the final frame is the trail of the big particle's centre over the preceding $60$ frames, showing its irregular, collision-driven path.}
\label{fig:task2}
\end{figure}

\subsection{Validation and limitations}
Unlike Task~1, the supplied program does not compute a mean-squared-displacement curve for the big particle, so no direct quantitative check of diffusive $\langle R^2\rangle \propto t$ scaling analogous to Eq.~\eqref{eq:msd} is available from its output; the only check performed by the program itself is the qualitative visual one described above. A natural extension -- not present in the submitted code -- would be to log the big particle's position at every frame and fit $\langle |\vec r(t)-\vec r(0)|^2\rangle$ against $t$, exactly as is done analytically for Task~1, to test whether the discrete collision dynamics reproduces the same linear growth law in the long-time limit; this is noted here as a limitation of the model as implemented rather than as a defect claimed to have been fixed.

Other idealisations are: perfectly elastic, instantaneous, pairwise collisions (no small-particle self-interactions, so the bath is not itself a genuine hard-sphere gas obeying the Maxwell--Boltzmann speed distribution -- every small particle keeps the same fixed speed $s$ throughout, only its direction is re-randomised); a phenomenological rather than Poissonian free-path model, as discussed above; and two-dimensional rather than three-dimensional confinement. These are reasonable simplifications for an illustrative animation but mean the simulation should be read as a qualitative, mechanistic illustration of the origin of Brownian motion rather than a quantitatively calibrated kinetic-theory calculation.

% =====================================================================
\section{Task 3 --- Planck Radiation and the Einstein Solid}
\label{sec:task3}
% =====================================================================

\subsection{Physical context}
Task~3 addresses two of the historically foundational results of early quantum theory: the Planck spectral distribution of black-body radiation, and the Einstein model of the specific heat of a solid. Both were introduced to resolve failures of classical (equipartition-based) statistical mechanics -- the ultraviolet catastrophe for radiation, and the failure of the Dulong--Petit law at low temperature for solids -- by quantising the energy of an oscillator in units of $\hbar\omega$.

\subsection{Theory: Planck spectrum}
For a single electromagnetic mode of angular frequency $\omega$ in thermal equilibrium at temperature $T$, quantising the mode as a simple harmonic oscillator with energy levels $E_n = n\hbar\omega$ gives a mean occupation number (the Bose--Einstein / Planck distribution)
\begin{equation}
\langle n\rangle = \frac{1}{e^{\hbar\omega/k_BT}-1}.
\end{equation}
Multiplying by the photon energy $\hbar\omega$ and the density of electromagnetic modes per unit volume and per unit angular frequency, $g(\omega)\,\dd\omega = \omega^2/(\pi^2 c^3)\,\dd\omega$ (two polarisations, summed over direction), gives the spectral energy density; converting to spectral radiance as a function of wavelength $\lambda = 2\pi c/\omega$ yields the standard Planck law
\begin{equation}
\label{eq:planck}
B(\lambda,T) = \frac{2hc^2}{\lambda^5}\,\frac{1}{e^{hc/\lambda k_B T}-1},
\end{equation}
with units of \si{\watt\per\square\meter\per\steradian\per\meter}. This is exactly the quantity implemented in \texttt{planck\_spectrum}.

\subsection{Theory: Einstein solid and its relation to the Debye model}
The Einstein model treats each of the $3N_A$ vibrational degrees of freedom of one mole of a monatomic solid as an independent quantum harmonic oscillator of the same frequency, characterised by an Einstein temperature $\theta_E = \hbar\omega_E/k_B$. The mean energy per oscillator is $\langle E\rangle = \hbar\omega_E/(e^{\hbar\omega_E/k_BT}-1)$, and differentiating the total energy $U=3N_A\langle E\rangle$ with respect to temperature gives the molar heat capacity
\begin{equation}
\label{eq:einstein_cv}
C_V(T) = 3R\left(\frac{\theta_E}{T}\right)^2 \frac{e^{-\theta_E/T}}{\left(1-e^{-\theta_E/T}\right)^2},
\end{equation}
which is exactly the expression evaluated by \texttt{einstein\_heat\_capacity} (rewritten there in terms of $e^{-x}$, $x=\theta_E/T$, for numerical stability at low $T$, avoiding overflow in $e^{+x}$ as $T\to0$). In the high-temperature limit $\theta_E/T\ll1$, expanding Eq.~\eqref{eq:einstein_cv} in powers of $x=\theta_E/T$ gives $C_V \to 3R\left[1-\tfrac{1}{12}x^2+O(x^4)\right]$, recovering the classical Dulong--Petit result $C_V\to3R$ as $T\to\infty$, and correctly vanishing exponentially (rather than as a power law) as $T\to0$, which was historically the qualitative success of Einstein's model over the classical equipartition prediction.

The (more accurate) Debye model, which treats the solid's normal modes with a realistic density of states $g(\omega)\propto\omega^2$ up to a cutoff $\omega_D$ rather than as $3N_A$ identical oscillators, has the corresponding high-temperature expansion $C_V \to 3R[1 - \tfrac{1}{20}(\theta_D/T)^2 + O(T^{-4})]$, with a characteristic Debye temperature $\theta_D$ for each material \cite{kittel}. Rather than treat $\theta_E$ as an independently tabulated material constant, \texttt{task3.py} derives an \emph{effective} single-oscillator Einstein temperature from the standard tabulated Debye temperature by matching the leading-order high-temperature quantum correction of the two models:
\begin{equation}
3R\left[1-\frac{1}{12}\left(\frac{\theta_E}{T}\right)^2\right] = 3R\left[1-\frac{1}{20}\left(\frac{\theta_D}{T}\right)^2\right]
\quad\Longrightarrow\quad
\theta_E = \sqrt{\frac{3}{5}}\,\theta_D .
\label{eq:thetaE}
\end{equation}
The Debye temperatures used, $\theta_D=\SI{170}{K}$ (Au), $\SI{343}{K}$ (Cu), $\SI{470}{K}$ (Fe), are standard tabulated solid-state values \cite{kittel}, giving effective Einstein temperatures $\theta_E = 131.7$, $265.7$ and $\SI{364.1}{K}$ respectively (reproduced directly from the legend of Fig.~\ref{fig:task3b}, where they are printed to one decimal place by the program). This is a physically motivated but explicitly approximate substitution: it guarantees the correct high-temperature ($T\gg\theta$) behaviour of $C_V(T)$ by construction, but the Einstein and Debye models have different low-temperature power laws ($C_V^{\rm Einstein}\sim e^{-\theta_E/T}$ versus the correct Debye $T^3$ law), so this matching procedure should not be expected to reproduce the correct low-temperature heat capacity; it is a deliberate simplification used only to obtain a physically reasonable characteristic temperature scale for each metal within the strictly single-frequency Einstein framework that the task specifies.

\subsection{Computational method and results}
Both spectra are evaluated on dense linearly spaced arrays (wavelength from $1$ to $\SI{3000}{nm}$ for the Planck spectrum, temperature from $5$ to $\SI{800}{K}$ for the heat capacities) using vectorised \texttt{NumPy} expressions, with the exponent clipped at $700$ (\texttt{np.expm1}, \texttt{np.exp}) to avoid floating-point overflow at short wavelength or low temperature -- a purely numerical safeguard that does not alter the mathematics in the regime where the exponent is $\lesssim700$ (double-precision \texttt{float64} already overflows for $e^x$ at $x\gtrsim709$, so the clip is inactive except in the extreme wing of the distributions and prevents a runtime warning without visibly truncating the plotted curves).

Figure~\ref{fig:task3a} shows the resulting Planck spectra for $T=4000$, $5000$ and $\SI{6000}{K}$. The peak wavelength shifts to shorter wavelength and the peak radiance increases sharply as $T$ increases, consistent with Wien's displacement law $\lambda_\text{peak}T \approx \SI{2.898e-3}{m.K}$ (not separately computed by the program, but visually consistent with the peaks near $\sim\SI{600}{nm}$--$\SI{700}{nm}$ at these temperatures) and with the Stefan--Boltzmann $T^4$ growth of the total radiated power.

\begin{figure}[H]
\centering
\includegraphics[width=0.85\linewidth]{figures/task3_planck_spectrum.png}
\caption{Planck spectral radiance $B(\lambda,T)$, Eq.~\eqref{eq:planck}, for three source temperatures, evaluated by \texttt{task3.py}.}
\label{fig:task3a}
\end{figure}

Figure~\ref{fig:task3b} shows the Einstein molar heat capacity, Eq.~\eqref{eq:einstein_cv}, for gold, copper and iron, using the Debye-derived Einstein temperatures of Eq.~\eqref{eq:thetaE}. All three curves rise from zero at low temperature and saturate at the Dulong--Petit value $3R = \SI{24.94}{J.mol^{-1}.K^{-1}}$ (grey dashed line) at high temperature, as required. The metal with the lowest characteristic temperature (gold, $\theta_E=\SI{131.7}{K}$) saturates at the lowest temperature, and iron (the stiffest of the three, $\theta_E=\SI{364.1}{K}$) saturates last -- correctly reproducing the qualitative trend that materials with stronger interatomic bonding (higher Debye/Einstein temperature) retain a lower heat capacity to higher temperature.

\begin{figure}[H]
\centering
\includegraphics[width=0.85\linewidth]{figures/task3_einstein_heat_capacity.png}
\caption{Einstein-model molar heat capacity $C_V(T)$, Eq.~\eqref{eq:einstein_cv}, for gold, copper and iron, using Einstein temperatures $\theta_E=\sqrt{3/5}\,\theta_D$ (Eq.~\eqref{eq:thetaE}) derived from standard tabulated Debye temperatures \cite{kittel}. The grey dashed line is the classical Dulong--Petit limit $3R$.}
\label{fig:task3b}
\end{figure}

\subsection{Limitations}
The Planck-spectrum calculation is exact within the stated (non-relativistic, thermal-equilibrium blackbody) assumptions, and no approximation beyond floating-point overflow protection is made. The Einstein-model calculation is exact for a genuine single-frequency Einstein solid, but as discussed above, the specific $\theta_E$ values used are not independently measured Einstein temperatures; they are derived from Debye temperatures via a high-temperature-matching approximation, Eq.~\eqref{eq:thetaE}, and the resulting curves should not be expected to reproduce the correct $T^3$ low-temperature behaviour of a real solid.

% =====================================================================
\section{Task 4 --- The Photoelectric Effect}
\label{sec:task4}
% =====================================================================

\subsection{Physical context and theory}
In the photoelectric effect, electrons are ejected from a metal surface when illuminated by light above a threshold frequency, with the maximum photoelectron kinetic energy $K_\text{max}$ increasing linearly with the frequency of the incident light, independent of intensity -- a result inexplicable in the classical (continuous-energy) wave picture of light, and explained by Einstein in 1905 by positing that light is absorbed in discrete quanta of energy $hf$ \cite{einstein1905photoelectric}. Einstein's photoelectric equation is
\begin{equation}
hf = \phi + K_\text{max},
\end{equation}
where $\phi$ is the work function of the metal (the minimum energy needed to remove an electron from the surface). Photoelectrons are conventionally characterised experimentally by the stopping voltage $V_s$, the retarding potential just sufficient to prevent any photoelectron reaching the collector, defined by $K_\text{max}=eV_s$. Combining these,
\begin{equation}
\label{eq:stopping}
V_s = \frac{h}{e}f - \frac{\phi}{e} = \frac{h}{e}f - \phi_\eV,
\end{equation}
where $\phi_\eV$ is the work function expressed directly in electron-volts (so that $\phi/e$ has the numerical value $\phi_\eV$ in volts). Below the threshold frequency $f_0=\phi/h$, no photoelectrons are emitted regardless of intensity, and $V_s$ is undefined (rather than negative).

\subsection{Computational method}
\texttt{stopping\_voltage} in \texttt{task4.py} evaluates Eq.~\eqref{eq:stopping} on an array of incident frequencies for five metals with standard tabulated work functions \cite{hyperphysics}: caesium ($\SI{2.10}{eV}$), sodium ($\SI{2.28}{eV}$), zinc ($\SI{4.30}{eV}$), copper ($\SI{4.70}{eV}$) and platinum ($\SI{6.35}{eV}$). Frequencies below each metal's threshold $f_0=\phi/h$ are explicitly masked to \texttt{NaN} using \texttt{numpy.where}, so that \texttt{Matplotlib} simply omits them from the plotted line rather than drawing an unphysical negative stopping voltage.

\subsection{Results}
Figure~\ref{fig:task4} shows the five stopping-voltage lines. All five are parallel straight lines of gradient $h/e = \SI{4.136e-15}{V.s}$ (Eq.~\eqref{eq:stopping}), each truncated at its own threshold frequency and each intersecting $V_s=0$ at $f=f_0=\phi/h$ -- the qualitative signature (common slope $h/e$, metal-dependent intercept) that historically allowed Millikan to measure $h$ from exactly this kind of plot. The printed threshold frequencies,
\begin{quote}\ttfamily
Cesium (Cs): 5.0778e+14 Hz \quad
Sodium (Na): 5.5130e+14 Hz \quad
Zinc (Zn): 1.0397e+15 Hz\\
Copper (Cu): 1.1365e+15 Hz \quad
Platinum (Pt): 1.5354e+15 Hz
\end{quote}
increase monotonically with $\phi$, exactly as $f_0=\phi/h$ requires, and can be checked directly: for caesium, $f_0 = (2.10\times1.602176634\times10^{-19})/(6.62607015\times10^{-34}) = \SI{5.078e14}{Hz}$, matching the printed value to the quoted precision.

\begin{figure}[H]
\centering
\includegraphics[width=0.85\linewidth]{figures/task4_photoelectric_effect.png}
\caption{Photoelectron stopping voltage $V_s$ against incident photon frequency for five metals, Eq.~\eqref{eq:stopping}. Each line begins at its own threshold frequency $f_0=\phi/h$ and has the same slope $h/e$.}
\label{fig:task4}
\end{figure}

An interactive extension of this task, providing a PhET-style animated model in which the metal, photon frequency, intensity and applied stopping voltage can all be varied by the user, is described in Section~\ref{sec:ext4}.

\subsection{Limitations}
The model assumes a single sharp work function per metal (real surfaces have a distribution of work functions depending on crystal face and contamination), non-relativistic photoelectron kinematics, and ignores the frequency-dependence of photon absorption probability (quantum yield) -- the plot shows only the kinematic upper bound $K_\text{max}$, not the photocurrent magnitude, which does depend on intensity in the way that $V_s$ itself does not.

% =====================================================================
\section{Task 5 --- The Bohr Model and the Hydrogen Emission Spectrum}
\label{sec:task5}
% =====================================================================

\subsection{Physical context}
Bohr's 1913 model of the hydrogen atom \cite{bohr1913} combines classical circular-orbit mechanics with an ad hoc quantisation condition on angular momentum, predicting a discrete set of stable electron orbits and, correctly, the observed line spectrum of atomic hydrogen. It is not a correct quantum-mechanical treatment (the true ground state has zero orbital angular momentum, not $\hbar$), but it correctly predicts the hydrogen energy levels because of a fortunate cancellation of errors, and remains the standard first derivation of the Rydberg formula.

\subsection{Derivation}
For an electron of mass $m_e$ and charge $-e$ in a circular orbit of radius $r$ about a proton, Newton's second law with the Coulomb attraction as the centripetal force gives
\begin{equation}
\frac{m_e v^2}{r} = \frac{e^2}{4\pi\varepsilon_0 r^2}. \label{eq:coulomb}
\end{equation}
Bohr's quantisation postulate is that orbital angular momentum is an integer multiple of $\hbar=h/2\pi$:
\begin{equation}
m_e v r = n\hbar, \qquad n=1,2,3,\ldots \label{eq:quantisation}
\end{equation}
Solving Eq.~\eqref{eq:quantisation} for $v=n\hbar/(m_e r)$ and substituting into Eq.~\eqref{eq:coulomb} gives
\begin{equation}
\frac{m_e}{r}\left(\frac{n\hbar}{m_e r}\right)^2 = \frac{e^2}{4\pi\varepsilon_0 r^2}
\quad\Longrightarrow\quad
r_n = \frac{4\pi\varepsilon_0\hbar^2 n^2}{m_e e^2}. \label{eq:rn}
\end{equation}
This is exactly the expression implemented in \texttt{bohr\_radius}. The total mechanical energy is $E=\tfrac12 m_ev^2 - e^2/(4\pi\varepsilon_0 r)$; using Eq.~\eqref{eq:coulomb} to write $\tfrac12 m_ev^2 = e^2/(8\pi\varepsilon_0 r)$ gives the standard virial-theorem result $E=-e^2/(8\pi\varepsilon_0 r)$, and substituting the quantised radius $r_n$ of Eq.~\eqref{eq:rn} gives
\begin{equation}
E_n = -\frac{m_e e^4}{8\varepsilon_0^2 h^2}\,\frac{1}{n^2} \equiv \frac{E_1}{n^2}. \label{eq:En}
\end{equation}
\texttt{task5.py} evaluates the ground-state energy $E_1$ directly from Eq.~\eqref{eq:En} in SI units, rather than inserting the textbook value $\SI{13.6}{eV}$ by hand; a downward transition from level $n_i$ to level $n_f<n_i$ then emits a photon of energy $E_\gamma = E_{n_i}-E_{n_f} = h c/\lambda$, giving the emitted wavelength
\begin{equation}
\frac{1}{\lambda} = \frac{m_ee^4}{8\varepsilon_0^2h^3c}\left(\frac{1}{n_f^2}-\frac{1}{n_i^2}\right),
\end{equation}
the Rydberg formula, with the Rydberg constant $R_\infty = m_ee^4/(8\varepsilon_0^2h^3c)$ appearing automatically as the combination of constants multiplying the bracket, rather than being separately hard-coded.

\subsection{Computational method}
The ground-state energy is computed once, in SI units, from Eq.~\eqref{eq:En} using the CODATA values of $m_e$, $e$, $\varepsilon_0$ and $h$, and converted to electron-volts by dividing by $e$. All fifteen downward transitions among levels $n=2,\ldots,6$ (i.e., every pair $n_i>n_f$ with $2\le n_i\le6$, $1\le n_f<n_i$) are enumerated by a nested loop, sorted by increasing wavelength, and plotted as a scatter of photon energy against wavelength (log scale, since the transition wavelengths span more than two orders of magnitude, from the Lyman to the Brackett-and-beyond series). A second figure plots the first four energy levels as horizontal lines together with two representative transition arrows.

\subsection{Results and validation}
The printed ground-state energy and Bohr radius are
\begin{quote}\ttfamily
Calculated Bohr ground-state energy: -13.6056931 eV\\
Bohr radius: 0.529177 Angstrom
\end{quote}
These match the standard textbook values of $-\SI{13.6057}{eV}$ and $a_0=\num{0.529177}\,\text{\AA}$ to the precision quoted, and are consistent with (though numerically slightly different from) the measured hydrogen ionisation energy of $\SI{13.5984}{eV}$ \cite{nist_hydrogen}; the $\sim0.05\%$ discrepancy is the well-known effect of the Bohr model's use of the electron mass in place of the electron--proton reduced mass $\mu=m_em_p/(m_e+m_p)\approx0.99946\,m_e$ (a correction of the correct sign and order of magnitude, though not separately implemented in this program) together with the model's neglect of relativistic and quantum-electrodynamic corrections.

Table~\ref{tab:transitions} reproduces the full printed transition table (photon energy and vacuum wavelength for every $n_i\to n_f$ pair with $2\le n_i\le6$). The five transitions terminating on $n_f=1$ (the Lyman series) lie in the ultraviolet ($93$--$\SI{122}{nm}$); those terminating on $n_f=2$ (Balmer series) span the near-ultraviolet to red, including the well-known $3\to2$ transition at $\SI{656.1}{nm}$ (H$\alpha$) and $2\to1$ at $\SI{121.5}{nm}$ (Lyman-$\alpha$), both annotated in Fig.~\ref{fig:task5b}; the remaining transitions onto $n_f=3,4$ fall in the infrared, consistent with the Paschen and Brackett series.

\begin{table}[H]
\centering
\caption{Hydrogen emission transitions computed by \texttt{task5.py}, sorted by wavelength (as printed by the program).}
\label{tab:transitions}
\begin{tabular}{ccc}
\toprule
Transition & Wavelength / nm & Photon energy / eV \\
\midrule
$6\to1$ & 93.730 & 13.2278 \\
$5\to1$ & 94.924 & 13.0615 \\
$4\to1$ & 97.202 & 12.7553 \\
$3\to1$ & 102.518 & 12.0939 \\
$2\to1$ & 121.502 & 10.2043 \\
$6\to2$ & 410.070 & 3.0235 \\
$5\to2$ & 433.937 & 2.8572 \\
$4\to2$ & 486.009 & 2.5511 \\
$3\to2$ & 656.112 & 1.8897 \\
$6\to3$ & 1093.520 & 1.1338 \\
$5\to3$ & 1281.469 & 0.9675 \\
$4\to3$ & 1874.607 & 0.6614 \\
$6\to4$ & 2624.449 & 0.4724 \\
$5\to4$ & 4050.076 & 0.3061 \\
$6\to5$ & 7455.821 & 0.1663 \\
\bottomrule
\end{tabular}
\end{table}

\begin{figure}[H]
\centering
\begin{subfigure}[b]{0.62\linewidth}
\includegraphics[width=\linewidth]{figures/task5_hydrogen_emission_spectrum.png}
\caption{}
\label{fig:task5a}
\end{subfigure}
\hfill
\begin{subfigure}[b]{0.36\linewidth}
\includegraphics[width=\linewidth]{figures/task5_bohr_model.png}
\caption{}
\label{fig:task5b}
\end{subfigure}
\caption{(a) Hydrogen emission spectrum: photon energy against wavelength (log scale) for all $15$ transitions of Table~\ref{tab:transitions}, with Balmer-series ($n_f=2$) transitions labelled. (b) Bohr energy-level diagram for $n=1$--$4$, with the $3\to2$ (H$\alpha$, $\SI{656.5}{nm}$) and $2\to1$ (Lyman-$\alpha$, $\SI{121.6}{nm}$) transitions marked.}
\label{fig:task5}
\end{figure}

\subsection{Limitations}
The Bohr model is not the correct quantum-mechanical description of hydrogen (it predicts, for example, that the ground state has orbital angular momentum $\hbar$ rather than the correct value $0$), and the derivation above uses the electron mass rather than the reduced mass, and omits fine structure, the Lamb shift, and relativistic corrections; nonetheless, because Eq.~\eqref{eq:En} happens to reproduce the correct non-relativistic energy-level spacing of hydrogen (a consequence of the hydrogen atom's accidental $SO(4)$ degeneracy, not derivable from the Bohr model itself), the computed transition energies agree with the observed hydrogen spectrum to the percent level shown above.

% =====================================================================
\section{Task 6 --- Electron Diffraction from Polycrystalline Graphite}
\label{sec:task6}
% =====================================================================

\subsection{Physical context}
Accelerating electrons through a potential difference $V$ and diffracting them from a polycrystalline target directly demonstrates the de Broglie wave nature of matter, and is the basis of the classic graphite electron-diffraction tube experiment. Task~6 models this apparatus: electrons accelerated through $1$--$\SI{5}{kV}$, diffracted by the two characteristic interplanar spacings of graphite, and detected as rings on a spherical phosphor screen.

\subsection{Derivation}
An electron accelerated from rest through potential difference $V$ gains kinetic energy $eV = p^2/(2m_e)$, so its momentum is $p=\sqrt{2m_eeV}$ and its de Broglie wavelength is
\begin{equation}
\lambda = \frac{h}{p} = \frac{h}{\sqrt{2m_eeV}}. \label{eq:debroglie}
\end{equation}
Constructive interference from a family of crystal planes of spacing $d$ requires Bragg's law $n\lambda=2d\sin\phi$, where $\phi$ is the glancing (Bragg) angle and $n$ the diffraction order, so that
\begin{equation}
\sin\phi = \frac{n\lambda}{2d}. \label{eq:bragg}
\end{equation}
On a hemispherical phosphor screen of radius $r_\text{screen}=\SI{65}{mm}$ centred on the sample, a beam diffracted through total angle $2\phi$ from the straight-through direction strikes the screen at transverse distance $x=r_\text{screen}\sin(2\phi)$ from the centre, which is the ring radius plotted by \texttt{draw\_screen}. Because graphite is polycrystalline (a powder of randomly oriented microcrystallites), every possible orientation of the diffracting planes relative to the beam is present simultaneously, so the diffraction pattern is a set of concentric rings (one per allowed order $n$ and plane spacing $d$) rather than discrete spots, as it would be for a single crystal.

Combining Eqs.~\eqref{eq:debroglie} and \eqref{eq:bragg} to eliminate $\lambda$,
\begin{equation}
\sin\phi = \frac{n}{2d}\cdot\frac{h}{\sqrt{2m_eeV}}
\quad\Longrightarrow\quad
\frac{1}{\sqrt V} = \underbrace{\frac{2d\sqrt{2m_ee}}{nh}}_{\text{gradient}}\sin\phi, \label{eq:task6a}
\end{equation}
which is the ``Task~6a'' straight-line relationship implemented by the validation panel of the program: for a fixed diffraction order $n$ (here $n=1$), a plot of $1/\sqrt V$ against $\sin\phi$ is linear, with a gradient from which the plane spacing $d$ can be recovered by inverting Eq.~\eqref{eq:task6a}, exactly as \texttt{recovered\_spacing} does.

\subsection{Computational method}
For each of graphite's two characteristic spacings, $d_1=\SI{0.123}{nm}$ and $d_2=\SI{0.213}{nm}$ \cite{graphite_spacing}, \texttt{diffraction\_rings} enumerates \emph{every} diffraction order $n$ for which Eq.~\eqref{eq:bragg} has a solution -- that is, every integer $n$ up to $\lfloor 2d/\lambda\rfloor$, the point beyond which $\sin\phi>1$ and no real Bragg angle exists -- rather than truncating the order arbitrarily. At $\SI{1}{kV}$, for example, the wavelength is short enough that many diffraction orders are geometrically allowed for both spacings. The three screen panels are generated at $1$, $3$ and $\SI{5}{kV}$, plotting a ring at radius $r_\text{screen}\sin(2\phi)$ for every allowed $(n,d)$ pair; the fourth panel performs the Task~6a linear fit at fifteen voltages between $1$ and $\SI{5}{kV}$ using \texttt{numpy.polyfit}, for the first order ($n=1$) of each spacing separately.

\subsection{Results and validation}
Figure~\ref{fig:task6} shows the four-panel output. The linear fits recover $d=\SI{0.123}{nm}$ and $d=\SI{0.213}{nm}$ (legend of the validation panel) to the three decimal places quoted -- i.e. the fitted spacings reproduce the input spacings used to generate the synthetic $(\sin\phi, 1/\sqrt V)$ data essentially exactly, as expected for a noiseless least-squares fit to data generated from the same closed-form relationship, Eq.~\eqref{eq:task6a}, that the fit assumes. This confirms that the implementation of Eqs.~\eqref{eq:debroglie}--\eqref{eq:task6a} is self-consistent. The printed ring counts,
\begin{quote}\ttfamily
1 kV: 6 rings for d1; 10 rings for d2\\
3 kV: 10 rings for d1; 19 rings for d2\\
5 kV: 14 rings for d1; 24 rings for d2
\end{quote}
increase with accelerating voltage, as expected: higher $V$ gives shorter $\lambda$ (Eq.~\eqref{eq:debroglie}), which by Eq.~\eqref{eq:bragg} permits more orders $n$ to satisfy $\sin\phi\le1$ before the maximum order is reached.

\begin{figure}[H]
\centering
\includegraphics[width=0.92\linewidth]{figures/task6_electron_diffraction.png}
\caption{Modelled graphite electron-diffraction screens at $1$, $3$ and $\SI{5}{kV}$ (top left, top right, bottom left), showing every diffraction ring allowed by Eq.~\eqref{eq:bragg} for the two graphite spacings $d_1=\SI{0.123}{nm}$ (blue) and $d_2=\SI{0.213}{nm}$ (red). Bottom right: the Task~6a linear check of Eq.~\eqref{eq:task6a}, $1/\sqrt V$ against $\sin\phi$ for the first-order ring of each spacing, with least-squares fits recovering both input spacings to three decimal places.}
\label{fig:task6}
\end{figure}

\subsection{Limitations}
Because every geometrically allowed order is drawn with equal line weight, the modelled screens in Fig.~\ref{fig:task6} show considerably more rings, particularly at higher voltage, than would be visible in a real diffraction pattern: the code does not model the structure-factor and Debye--Waller intensity falloff that strongly suppresses high-order reflections in a real crystal, so the outer, high-order rings shown here would in practice be far too faint to observe. The model is also entirely non-relativistic (using $eV=p^2/2m_e$ rather than the relativistically correct electron energy-momentum relation), which introduces a growing error at higher accelerating voltage; at $\SI{5}{kV}$ the electron kinetic energy is small compared with its rest energy $m_ec^2\approx\SI{511}{keV}$ (a ratio of $\sim1\%$), so the non-relativistic approximation is reasonable at the voltages actually used here.

% =====================================================================
\section{Task 7 --- The Particle in a Box}
\label{sec:task7}
% =====================================================================

\subsection{Physical context and theory}
The infinite square well (particle in a box) is the simplest exactly solvable time-independent Schr\"odinger problem and the standard first example of energy quantisation, boundary-condition-enforced normalisation, and orthogonality of eigenstates in introductory quantum mechanics \cite{griffiths}. For a particle of mass $m$ confined to $0<x<a$ with infinite potential walls, the time-independent Schr\"odinger equation $-\frac{\hbar^2}{2m}\psi''(x)=E\psi(x)$ subject to $\psi(0)=\psi(a)=0$ has normalised solutions
\begin{equation}
\psi_n(x) = \sqrt{\frac2a}\sin\!\left(\frac{n\pi x}{a}\right), \qquad E_n = \frac{n^2\pi^2\hbar^2}{2ma^2}, \qquad n=1,2,3,\ldots \label{eq:pib}
\end{equation}
implemented directly by \texttt{energy} and \texttt{probability\_density} in \texttt{task7.py}, for an electron ($m=m_e$) confined to a $a=\SI{1}{nm}$ box.

\subsection{Computational method and results}
Both quantities in Eq.~\eqref{eq:pib} are evaluated on arrays: energies for $n=1,\ldots,6$ and probability densities $|\psi_n(x)|^2$ for $n=1,\ldots,4$ on a $1200$-point grid across the box (with zero density enforced outside $0<x<a$ by explicit masking). Figure~\ref{fig:task7} shows the resulting energy ladder (left) and probability densities (right).

\begin{figure}[H]
\centering
\includegraphics[width=\linewidth]{figures/task7_particle_in_a_box.png}
\caption{Left: particle-in-a-box energy levels $E_n\propto n^2$ for an electron in a $\SI{1}{nm}$ box, Eq.~\eqref{eq:pib}. Right: normalised probability densities $|\psi_n(x)|^2$ for the first four stationary states.}
\label{fig:task7}
\end{figure}

The plotted energies visibly follow the $n^2$ scaling of Eq.~\eqref{eq:pib}: reading from Fig.~\ref{fig:task7}, $E_1\approx\SI{0.38}{eV}$ and $E_4\approx\SI{6.0}{eV}\approx16E_1$, consistent with $E_4/E_1=4^2=16$. The probability densities show the expected number of nodes ($n-1$ internal zeros for state $n$, visible directly in the right-hand panel) and each integrates to unity over the box by construction (the normalisation constant $\sqrt{2/a}$ is built into Eq.~\eqref{eq:pib} rather than imposed numerically). An extension of this task, cross-checking the Heisenberg uncertainty principle for these same eigenstates both analytically and by direct numerical integration, is described in Section~\ref{sec:ext7}.

\subsection{Limitations}
The infinite-well potential is an idealisation (no real confining potential has infinite walls); the model also ignores electron spin and any interaction between the confined particle and its environment, and treats the problem non-relativistically, appropriate for the low energies ($\lesssim\SI{15}{eV}$) obtained here.

% =====================================================================
\section{Task 8 --- Photon Polarisation Measurement and Quantum Cryptography}
\label{sec:task8}
% =====================================================================

\subsection{Physical context}
Task~8 compares two competing predictions for the outcome of measuring the polarisation of a single photon, prepared in a definite linear polarisation state at angle $\theta$, using a polarising analyser oriented at angle $\phi$: a purely quantum-mechanical prediction (Malus's law, applied to a single-photon probability rather than a classical beam intensity), and a naive classical (hidden-variable) prediction in which the photon is imagined to carry a definite, pre-existing $H$/$V$ identity rather than existing in a coherent superposition. This contrast between a hidden-variable-type classical prediction and the correct quantum prediction for a mismatched measurement basis is directly relevant to quantum key distribution protocols such as BB84, in which the error rate introduced by measuring a photon in the wrong basis is exactly this mismatch probability \cite{nielsen_chuang}.

\subsection{Derivation}
Quantum mechanically, a photon prepared in the linear-polarisation state $\ket{\theta}=\cos\theta\ket{H}+\sin\theta\ket{V}$ measured by a polariser/detector pair aligned with $\ket{\phi}=\cos\phi\ket{H}+\sin\phi\ket{V}$ is transmitted (detected as ``$\phi$'') with the Born-rule probability
\begin{equation}
P_\text{trans}^\text{quantum} = |\langle\phi|\theta\rangle|^2 = \cos^2(\phi-\theta),
\end{equation}
Malus's law applied at the single-photon level; the complementary ``mismatch'' probability (the photon registers in the orthogonal channel) is therefore
\begin{equation}
P_\text{mismatch}^\text{quantum} = \sin^2(\phi-\theta), \label{eq:qmismatch}
\end{equation}
exactly the expression implemented by \texttt{quantum\_mismatch}. A naive classical (hidden-variable) model instead supposes that the photon has already ``decided'', at the moment of preparation, to behave as either purely $H$- or purely $V$-polarised, with probabilities $\cos^2\theta$ and $\sin^2\theta$ respectively (an \emph{incoherent} classical mixture, rather than the coherent superposition $\ket\theta$ used above); each classical alternative is then independently subject to a Malus-law-type transmission probability through the analyser, $\cos^2\phi$ for the $H$ alternative and $\sin^2\phi$ for the $V$ alternative. The resulting classical ``match'' probability is
\begin{equation}
P_\text{match}^\text{classical} = \cos^2\theta\cos^2\phi+\sin^2\theta\sin^2\phi
\quad\Longrightarrow\quad
P_\text{mismatch}^\text{classical} = 1-\cos^2\theta\cos^2\phi-\sin^2\theta\sin^2\phi, \label{eq:cmismatch}
\end{equation}
matching \texttt{classical\_mismatch}. Comparing Eqs.~\eqref{eq:qmismatch} and \eqref{eq:cmismatch} directly is instructive: expanding the quantum match probability,
\begin{equation}
1-P_\text{mismatch}^\text{quantum} = \cos^2(\phi-\theta) = \cos^2\theta\cos^2\phi + \sin^2\theta\sin^2\phi + \underbrace{2\sin\theta\cos\theta\sin\phi\cos\phi}_{\tfrac12\sin2\theta\sin2\phi},
\end{equation}
shows that the quantum match probability equals the classical match probability of Eq.~\eqref{eq:cmismatch} \emph{plus} an interference (coherence) cross-term $\tfrac12\sin2\theta\sin2\phi$ that is entirely absent from the incoherent hidden-variable picture. The two predictions therefore coincide only when $\theta$ or $\phi$ is a multiple of $45^\circ$ (where the cross-term vanishes) and otherwise genuinely differ, which is exactly the discrepancy the task's calculator is designed to expose.

\subsection{Computational method and results}
The program plots both mismatch probabilities, Eqs.~\eqref{eq:qmismatch}--\eqref{eq:cmismatch}, as continuous functions of Bob's analyser angle $\phi$ for Alice's angle fixed at $\theta=-30^\circ$ (left panel of Fig.~\ref{fig:task8}), and additionally provides a two-bar ``calculator'' (right panel) evaluating both formulas at a specific, user-editable pair of angles ($\theta=-30^\circ$, $\phi=30^\circ$ in the version run here). The printed/plotted values,
\begin{quote}\ttfamily
Classical mismatch = 0.375 \qquad Quantum mismatch = 0.750,
\end{quote}
can be checked directly from Eqs.~\eqref{eq:qmismatch}--\eqref{eq:cmismatch}: with $\phi-\theta=60^\circ$, $\sin^2(60^\circ)=3/4=0.750$, matching the quantum bar exactly; and $\cos^2(-30^\circ)\cos^2(30^\circ)+\sin^2(-30^\circ)\sin^2(30^\circ) = (3/4)(3/4)+(1/4)(1/4)=9/16+1/16=10/16=0.625$, so $1-0.625=0.375$, matching the classical bar exactly. At this particular angle pair the two predictions differ by a factor of two, illustrating that the choice of measurement basis mismatch can produce a large, experimentally distinguishable difference between the classical hidden-variable prediction and the quantum prediction.

\begin{figure}[H]
\centering
\includegraphics[width=\linewidth]{figures/task8_quantum_cryptography.png}
\caption{Left: classical (Eq.~\eqref{eq:cmismatch}) and quantum (Eq.~\eqref{eq:qmismatch}) mismatch probability against Bob's analyser angle $\phi$, for Alice's photon fixed at $\theta=-30^\circ$; markers show the specific pair evaluated in the right-hand panel. Right: bar-chart calculator for $\theta=-30^\circ$, $\phi=30^\circ$ ($\phi-\theta=60^\circ$), giving classical mismatch $0.375$ and quantum mismatch $0.750$.}
\label{fig:task8}
\end{figure}

\subsection{Limitations}
The classical model implemented here is a specific, simple hidden-variable model (an incoherent $H$/$V$ mixture with independent Malus-law transmission of each component); it is not the most general local-hidden-variable model conceivable, and the task does not implement a Bell-type (CHSH) inequality test with entangled photon \emph{pairs}, which is the historically decisive experimental discriminator between quantum mechanics and local hidden-variable theories \cite{aspect1982}. The comparison here is at the level of a single photon and a single measurement, illustrating the general principle (superposition versus classical mixture) relevant to basis-mismatch errors in quantum key distribution, rather than a full entanglement-based Bell test.

% =====================================================================
\section{Task 9 --- Relativistic Compton Scattering}
\label{sec:task9}
% =====================================================================

\subsection{Physical context}
Compton scattering -- the increase in wavelength of a photon scattered from a (nearly) free electron -- was, with the photoelectric effect, one of the decisive early confirmations that light carries momentum as well as energy in discrete quanta \cite{compton1923}. Task~9 computes the full relativistic kinematics of $\SI{100}{keV}$ photons Compton-scattered through all angles $0\le\theta\le180^\circ$.

\subsection{Derivation}
Conservation of four-momentum for a photon of initial energy $E_\gamma=hc/\lambda$ scattering off a stationary electron of rest energy $m_ec^2$ gives the Compton formula for the wavelength shift,
\begin{equation}
\Delta\lambda = \lambda' - \lambda = \frac{h}{m_ec}(1-\cos\theta), \label{eq:compton}
\end{equation}
where $\theta$ is the photon scattering angle and $\lambda_C\equiv h/(m_ec)=\SI{2.426}{pm}$ is the Compton wavelength, implemented directly as \texttt{compton\_wavelength} in the code. Given the scattered wavelength $\lambda'=\lambda+\Delta\lambda$ and hence the scattered photon energy $E_\gamma'=hc/\lambda'$, conservation of momentum in the plane of scattering (photon momentum $p_\gamma=E_\gamma/c$) resolves into components along ($x$) and perpendicular to ($y$) the incident direction:
\begin{equation}
p_x = \frac{E_\gamma}{c}-\frac{E_\gamma'}{c}\cos\theta, \qquad p_y = \frac{E_\gamma'}{c}\sin\theta,
\end{equation}
which is the recoil electron's momentum, and the recoil angle is $\phi_e=\arctan(p_y/p_x)$ (\texttt{np.arctan2} in the code, correctly handling all quadrants). By conservation of energy, the electron's kinetic energy is $K=E_\gamma-E_\gamma'$, so its relativistic Lorentz factor and speed follow from
\begin{equation}
\gamma = 1+\frac{K}{m_ec^2}, \qquad v = c\sqrt{1-\gamma^{-2}}, \label{eq:relativistic_speed}
\end{equation}
the standard relativistic energy--momentum relation $E=\gamma m_ec^2$ inverted for speed, exactly as implemented in \texttt{compton\_results}. The full relativistic treatment is essential here (rather than the non-relativistic approximation used for the electron in Tasks~5--7): at $\SI{100}{keV}$ the incident photon energy is already $\sim20\%$ of the electron rest energy $m_ec^2=\SI{511}{keV}$, and, as Fig.~\ref{fig:task9} shows, the recoil electron speed approaches $v/c\approx0.32$ at backscattering, comfortably in the regime where $\gamma-1$ is not small and a non-relativistic treatment would be inaccurate.

\subsection{Computational method and results}
All quantities are evaluated as vectorised functions of the scattering angle $\theta$ on a $1000$-point grid from $0$ to $\pi$. Figure~\ref{fig:task9} shows the three resulting curves. The fractional wavelength shift $\Delta\lambda/\lambda$ (left panel) rises monotonically from $0$ at forward scattering ($\theta=0$, no momentum transfer) to a maximum at $\theta=180^\circ$ (backscattering), where Eq.~\eqref{eq:compton} gives the largest possible shift, $\Delta\lambda=2\lambda_C$; because the incident wavelength $\lambda=hc/E_\gamma=\SI{12.4}{pm}$ at $\SI{100}{keV}$ is only a few times the Compton wavelength $\lambda_C=\SI{2.43}{pm}$, the fractional shift $\Delta\lambda/\lambda$ reaches a substantial fraction ($\approx0.39$, read from the figure) at backscattering -- unlike for optical photons, where $\lambda\gg\lambda_C$ and the fractional shift is negligible, which is why the Compton effect is observable only for X-ray or higher-energy photons. The electron recoil speed (centre panel) increases monotonically with scattering angle, saturating near $v/c\approx0.32$ at $\theta=180^\circ$ (the angle of maximum momentum and energy transfer to the electron), and the recoil angle (right panel) decreases monotonically from $90^\circ$ (a glancing photon scatter transfers momentum purely transverse to the beam) to $0^\circ$ (a photon backscattered straight back transfers momentum straight forward), as required by transverse-momentum conservation.

\begin{figure}[H]
\centering
\includegraphics[width=\linewidth]{figures/task9_compton_scattering.png}
\caption{Relativistic Compton scattering of a $\SI{100}{keV}$ photon: fractional wavelength shift (left, Eq.~\eqref{eq:compton}), recoil electron speed as a fraction of $c$ (centre, Eq.~\eqref{eq:relativistic_speed}), and recoil electron angle (right), all as functions of photon scattering angle $\theta$.}
\label{fig:task9}
\end{figure}

\subsection{Validation}
Two independent checks are available without further computation. First, at $\theta=0$ every curve correctly gives the null result expected for forward (undeviated) scattering: $\Delta\lambda/\lambda=0$, $v=0$, matching Fig.~\ref{fig:task9}. Second, the transverse-momentum-conservation identity $p_\gamma\sin\theta = p_e\sin\phi_e$ used implicitly in $\phi_e=\arctan(p_y/p_x)$ guarantees $\phi_e=90^\circ$ exactly at $\theta\to0^+$ (a small-angle photon deflection is entirely transverse, so the recoiling electron -- which must carry the transverse momentum the photon has just acquired, since the electron started at rest with no transverse momentum -- moves purely transversely too) and $\phi_e\to0^\circ$ as $\theta\to180^\circ$, both of which are visible in the right-hand panel of Fig.~\ref{fig:task9}.

\subsection{Limitations}
The scattering electron is assumed initially at rest and free (i.e.\ unbound), appropriate for a high-energy photon scattering from a weakly bound outer-shell electron but not for scattering from tightly bound inner-shell electrons, where the binding energy is a non-negligible fraction of the photon energy and Compton scattering competes with photoelectric absorption and coherent (Rayleigh/Thomson) scattering, none of which are modelled here. The calculation also treats the scattering as a single, elastic, two-body kinematic event and does not compute a scattering cross-section (the Klein--Nishina formula), so no information about the relative \emph{probability} of scattering into a given angle is provided by this program, only the kinematics conditional on a given scattering angle having occurred.

% =====================================================================
\section{Task 10 --- Hydrogenic Orbital Probability Densities}
\label{sec:task10}
% =====================================================================

\subsection{Physical context}
Task~10 evaluates the full three-dimensional quantum-mechanical probability density $|\psi_{nlm}(\vec r)|^2$ for hydrogenic (single-electron, $Z=1$) orbitals, going beyond the radially symmetric, angular-momentum-free Bohr picture of Task~5 to the genuine three-quantum-number ($n,l,m$) solutions of the time-independent Schr\"odinger equation in a Coulomb potential.

\subsection{Theory}
The hydrogen atom's stationary states separate in spherical polar coordinates as $\psi_{nlm}(r,\vartheta,\varphi) = R_{nl}(r)\,Y_l^m(\vartheta,\varphi)$, where the radial wavefunction involves an associated Laguerre polynomial and the angular part is a spherical harmonic \cite{griffiths}:
\begin{equation}
R_{nl}(r) = \sqrt{\left(\frac{2}{na_0}\right)^3\frac{(n-l-1)!}{2n(n+l)!}}\;e^{-\rho/2}\rho^l\,L_{n-l-1}^{2l+1}(\rho), \qquad \rho \equiv \frac{2r}{na_0}, \label{eq:radial}
\end{equation}
with $a_0=4\pi\varepsilon_0\hbar^2/(m_ee^2)$ the Bohr radius (the same quantity derived in Eq.~\eqref{eq:rn} for $n=1$), and
\begin{equation}
|Y_l^m(\vartheta,\varphi)|^2 = \frac{2l+1}{4\pi}\frac{(l-|m|)!}{(l+|m|)!}\left[P_l^{|m|}(\cos\vartheta)\right]^2, \label{eq:angular}
\end{equation}
independent of $\varphi$ because only $|\psi|^2$, not $\psi$ itself, is required. Equations~\eqref{eq:radial}--\eqref{eq:angular} are exactly the quantities computed by \texttt{probability\_density}: the generalised Laguerre polynomial $L_{n-l-1}^{2l+1}(\rho)$ is evaluated via the three-term recurrence relation
\begin{equation}
kL_k^\alpha(x) = (2k-1+\alpha-x)L_{k-1}^\alpha(x) - (k-1+\alpha)L_{k-2}^\alpha(x),
\end{equation}
and the associated Legendre function $P_l^{|m|}(\cos\vartheta)$ via the analogous recurrence
\begin{equation}
(l-m)P_l^m(x) = (2l-1)xP_{l-1}^m(x) - (l+m-1)P_{l-2}^m(x),
\end{equation}
seeded at $P_m^m(x)=(-1)^m(2m-1)!!(1-x^2)^{m/2}$ and $P_{m+1}^m(x)=x(2m+1)P_m^m(x)$. This is a self-contained implementation of both special-function families using only elementary recurrence, without \texttt{scipy.special} or any other external special-function library.

\subsection{Computational method}
\texttt{probability\_density} accepts Cartesian coordinates $(x,y,z)$ directly (converting internally to $r=\sqrt{x^2+y^2+z^2}$ and $\cos\vartheta=z/r$), which makes it straightforward to evaluate the density on a two-dimensional Cartesian grid (for the 2-D cross-section plots) or a three-dimensional Cartesian grid (for the 3-D scatter plots) without any coordinate-mesh transformation. Four orbitals are studied: $1s$ ($n{=}1,l{=}0,m{=}0$), $2p_z$ ($n{=}2,l{=}1,m{=}0$), $3d_{z^2}$ ($n{=}3,l{=}2,m{=}0$) and an $n=4,l=3,m=2$ ($4f$) state.

The 2-D figure evaluates $|\psi|^2$ on the $y=0$ ($x$--$z$) plane, deliberately chosen over the equatorial $z=0$ plane because it is the plane in which orbitals whose angular factor vanishes at the equator (such as $2p_z$) are still visible. The 3-D figure retains only grid points where the density exceeds $15\%$ of that orbital's maximum density, the visualisation threshold specified by the challenge brief, and colours the retained points by their (renormalised) density value.

\subsection{Results}
Figure~\ref{fig:task10a} shows the four $x$--$z$ cross-sections. The $1s$ state is a compact, spherically symmetric blob centred on the origin, as expected for $l=0$ (the angular factor $Y_0^0$ is a constant, so $|\psi_{100}|^2$ depends only on $r$). The $2p_z$ state shows the expected two-lobed structure along $z$ with a node in the $z=0$ plane (the angular factor $\propto\cos^2\vartheta$ vanishes at $\vartheta=90^\circ$). The panel labelled $3d_{z^2}$ ($n{=}3,l{=}2,m{=}0$) shows two bright lobes along $z$ separated by a dark ``waist''; this is the same $x$--$z$ cross-section through the polar lobes of the familiar $d_{z^2}$ orbital, and its faint equatorial torus does not disappear (the angular factor $\propto(3\cos^2\vartheta-1)^2$ is $1/4$ of its polar value at the equator, not zero, which was confirmed by direct evaluation of the density along the $x$-axis at fixed $r$: exactly one quarter of the on-axis value at the same radius, for every radius tested) but is visually suppressed by the linear colour scale used for the plot, which is normalised to the much brighter polar maximum. The $n{=}4,l{=}3,m{=}2$ state shows four lobes arranged around the origin, reflecting its more complex angular nodal structure.

\begin{figure}[H]
\centering
\includegraphics[width=0.85\linewidth]{figures/task10_hydrogenic_orbitals_2d.png}
\caption{Hydrogenic probability density $|\psi_{nlm}|^2/|\psi_{nlm}|^2_\text{max}$ on the $x$--$z$ ($y=0$) plane, for the four orbitals listed in the text. Distances are in units of the Bohr radius $a_0$.}
\label{fig:task10a}
\end{figure}

Figure~\ref{fig:task10b} shows the corresponding 3-D scatter plots for $2p_z$ and $3d_{z^2}$, retaining points above $15\%$ of the maximum density. The $2p_z$ plot shows a dumbbell-shaped distribution of points along the $z$-axis, consistent with the two lobes seen in the cross-section, though the individual points are visibly arranged on the coarse $45\times45\times45$ sampling grid used for the 3-D calculation (a much sparser grid than the $260\times260$ grid used for the 2-D cross-sections), so the lobes appear grainy and cube-textured rather than as a smooth surface. The $3d_{z^2}$ plot makes the equatorial torus, only inferred indirectly from the 2-D cross-section, directly visible in three dimensions as a faint ring of points connecting the two bright polar lobes -- since the $15\%$-of-maximum threshold is well below the equatorial density's $25\%$-of-polar-maximum value derived above, the torus survives the thresholding and appears explicitly in the 3-D rendering even though it is difficult to discern in the corresponding 2-D cross-section of Fig.~\ref{fig:task10a}.

\begin{figure}[H]
\centering
\includegraphics[width=\linewidth]{figures/task10_hydrogenic_orbitals_3d.png}
\caption{Three-dimensional probability-density point clouds for the $2p_z$ (left) and $3d_{z^2}$ (right) orbitals, retaining only grid points with density above $15\%$ of that orbital's maximum, coloured by (renormalised) density.}
\label{fig:task10b}
\end{figure}

\subsection{Validation}
As a direct numerical check of the angular-factor claim made above, evaluating \texttt{probability\_density} for the $3d_{z^2}$ state along the polar axis ($x=y=0$, varying $z$) and along the equatorial axis ($z=y=0$, varying $x$) at ten matched radii from $0.01\,a_0$ to $20\,a_0$ gives a ratio (equatorial density)/(polar density) of exactly $0.25$ at every radius tested, in agreement with the analytic angular-factor ratio $\left[P_2^0(0)\right]^2/\left[P_2^0(1)\right]^2 = (-1/2)^2/(1)^2 = 1/4$ derived from Eq.~\eqref{eq:angular}. This both confirms the correctness of the associated-Legendre recurrence implementation and explains quantitatively why the equatorial torus of $3d_{z^2}$ is faint (but present) in the 2-D cross-section and yet clearly visible in the 3-D thresholded plot, as discussed above.

\subsection{Limitations}
The 3-D visualisation uses a comparatively coarse $45^3$ spatial grid (chosen, reasonably, to keep the scatter plot's point count and render time manageable) and a hard $15\%$ density threshold rather than a continuous opacity or isosurface rendering, both of which are visualisation choices rather than physical approximations; the underlying wavefunction evaluation itself, Eqs.~\eqref{eq:radial}--\eqref{eq:angular}, is exact (to floating-point precision) for the non-relativistic hydrogenic Schr\"odinger equation, subject to the usual neglect of electron spin, fine structure, and (for $Z>1$ ions, not used here) electron--electron interaction.

% =====================================================================
\section{Extensions}
\label{sec:extensions}
% =====================================================================

\subsection{Task 4 extension: an interactive photoelectric-effect simulator}
\label{sec:ext4}
The Task~4 extension (\texttt{task4\_extension.py}) implements an animated, interactive desktop application, built with Python's \texttt{tkinter} GUI toolkit, modelled explicitly on PhET-style physics demonstrations. Unlike every other program in this collection, it produces no static output file; its purpose is real-time user interaction, so it is documented here rather than accompanied by a generated figure.

The application exposes four user controls: a dropdown selecting one of the same five metals used in Task~4 (with the same tabulated work functions \cite{hyperphysics}), and three sliders for incident photon frequency ($3\times10^{14}$--$2\times10^{15}\,\si{Hz}$), light intensity ($0$--$100$, arbitrary units), and an applied (retarding) stopping voltage ($0$--$\SI{10}{V}$). The underlying physics functions -- \texttt{threshold\_frequency}, \texttt{photon\_energy\_eV}, \texttt{maximum\_kinetic\_energy\_eV} and \texttt{stopping\_voltage} -- are the same Einstein-equation calculations, Eq.~\eqref{eq:stopping}, as in the non-interactive Task~4 program, re-implemented here to operate on scalar (rather than array) inputs suited to a live GUI callback.

The animation loop, driven by \texttt{tkinter}'s \texttt{after} scheduler at a fixed $\SI{30}{ms}$ cadence, spawns simulated photon markers from the light-source side of the canvas at a rate controlled by the intensity slider; when a photon marker reaches the metal-plate boundary, the program checks whether the current frequency exceeds the threshold frequency for the selected metal, and if so spawns an electron marker with an initial speed sampled (uniformly, between $15\%$ and $100\%$ of the maximum) to represent the physical spread of photoelectron kinetic energies up to $K_\text{max}=hf-\phi$; if the applied stopping voltage is at or above $K_\text{max}/e$, emitted electrons are decelerated (an exponential speed decay is applied per frame) rather than being allowed to reach the collector electrode, illustrating the stopping-voltage measurement principle directly. A text readout beneath the canvas continuously displays the work function, current photon energy, threshold frequency, and $K_\text{max}$, together with a status line indicating whether photoelectrons are currently being emitted, stopped, or entirely absent (below threshold).

As the program's own docstring makes clear, the sampled initial electron speed and the frame-by-frame deceleration under an applied stopping voltage are deliberately illustrative rather than physically first-principles: real photoelectron energies are not uniformly distributed up to $K_\text{max}$ (the distribution depends on the material's band structure and the depth from which the electron is emitted), and the deceleration model does not solve the electron's actual equation of motion under the retarding field. The value of the extension lies in giving an immediate, manipulable visual and numerical correspondence to Einstein's equation -- in particular, making directly visible the sharp threshold behaviour (no electrons at all below $f_0$, regardless of intensity) that is the qualitative signature distinguishing the photon picture from a classical wave picture, in which any sufficiently intense low-frequency light would eventually eject electrons.

\subsection{Task 7 extension: verifying the Heisenberg uncertainty principle}
\label{sec:ext7}
The Task~7 extension (\texttt{task7\_extension.py}) tests the Heisenberg uncertainty relation $\Delta x\,\Delta p\ge\hbar/2$ for the particle-in-a-box eigenstates of Eq.~\eqref{eq:pib}, computing the uncertainty product $\Delta x\,\Delta p$ two independent ways -- from a closed-form analytic expression, and from direct numerical integration of the wavefunction -- and comparing both against the Heisenberg bound.

\subsubsection{Analytic formula}
For the normalised state $\psi_n(x)=\sqrt{2/a}\sin(n\pi x/a)$ on $[0,a]$, the standard textbook calculation (see, e.g., \cite{griffiths}) gives $\langle x\rangle = a/2$ (by the symmetry of $|\psi_n|^2$ about the box centre) and
\begin{equation}
\langle x^2\rangle = \frac{a^2}{3} - \frac{a^2}{2n^2\pi^2},
\end{equation}
so that
\begin{equation}
(\Delta x)^2 = \langle x^2\rangle - \langle x\rangle^2 = \frac{a^2}{3}-\frac{a^2}{2n^2\pi^2} - \frac{a^2}{4} = a^2\left(\frac{1}{12}-\frac{1}{2n^2\pi^2}\right). \label{eq:correct_dx}
\end{equation}
With $\langle p\rangle=0$ (exact by symmetry for a real stationary state) and $\langle p^2\rangle = 2mE_n = (n\pi\hbar/a)^2$ from Eq.~\eqref{eq:pib}, so that $\Delta p = n\pi\hbar/a$, multiplying $\Delta x$ from Eq.~\eqref{eq:correct_dx} by $\Delta p$ gives the closed-form uncertainty product
\begin{equation}
\Delta x\,\Delta p = n\pi\hbar\sqrt{\frac{1}{12}-\frac{1}{2n^2\pi^2}} = \frac{n\pi\hbar}{2}\sqrt{\frac13-\frac{2}{\pi^2n^2}}, \label{eq:code_formula}
\end{equation}
which is exactly the expression implemented by \texttt{analytic\_uncertainty\_product}. At $n=1$ the term under the square root is $\tfrac13-\tfrac{2}{\pi^2}=0.1307>0$, so, unlike the naively simplified form $\tfrac13-\tfrac{1}{2n^2}$ that a careless rearrangement of Eq.~\eqref{eq:correct_dx} might produce (which would erroneously drop the factor of $\pi^2$ from the second term and go negative at $n=1$), Eq.~\eqref{eq:code_formula} is well-defined and real for every $n\ge1$: keeping the $\pi^2$ in the denominator of the second term is essential to obtaining a physically sensible ground-state result.

\subsubsection{Numerical cross-check}
Independently, \texttt{numerical\_uncertainty\_product} computes $\Delta x=\sqrt{\langle x^2\rangle-\langle x\rangle^2}$ by trapezoidal integration (\texttt{numpy.trapezoid}) of $x\,|\psi_n(x)|^2$ and $x^2|\psi_n(x)|^2$ over a $10^5$-point grid, and takes $\Delta p = n\pi\hbar/a$ directly, exactly as derived above. Because this branch integrates the wavefunction numerically rather than evaluating the closed-form expression, it provides a genuinely independent check on Eq.~\eqref{eq:code_formula}.

Running the program prints, for the ground state,
\begin{quote}\ttfamily
n=1: Delta-x Delta-p = 0.567862 hbar
\end{quote}
consistent with evaluating Eq.~\eqref{eq:code_formula} directly at $n=1$: $\hbar\sqrt{\pi^2/12-1/2}=0.5679\,\hbar$, comfortably above the Heisenberg bound $\hbar/2=0.5\,\hbar$.

\subsubsection{Results}
Figure~\ref{fig:task7ext} shows both series across $n=1,\ldots,12$. The analytic curve (blue circles) and the numerically integrated points (red crosses) coincide at every plotted $n$, including $n=1$, to within the marker size, confirming that the closed-form expression, Eq.~\eqref{eq:code_formula}, and the independent numerical integration of the exact wavefunction agree. Both series lie above the Heisenberg bound $\hbar/2$ (black dashed line) at every $n$, with the uncertainty product growing approximately linearly with $n$ for $n\gtrsim3$ (as expected from the $n\pi\hbar/2\cdot\sqrt{1/3}$ leading behaviour of Eq.~\eqref{eq:code_formula} once the $O(1/n^2)$ correction becomes small) and approaching the bound most closely at $n=1$, where the correction term is largest relative to the leading term.

\begin{figure}[H]
\centering
\includegraphics[width=0.8\linewidth]{figures/task7_uncertainty_principle.png}
\caption{Uncertainty product $\Delta x\,\Delta p/\hbar$ for particle-in-a-box eigenstates $n=1,\ldots,12$: the closed-form analytic formula, Eq.~\eqref{eq:code_formula} (blue circles), and an independent numerical integration of the wavefunction (red crosses), compared with the Heisenberg bound $\hbar/2$ (black dashed line). The two series coincide at every $n$, including the ground state.}
\label{fig:task7ext}
\end{figure}

The two independent methods agreeing at every $n$, including the ground state where the correction term in Eq.~\eqref{eq:code_formula} is largest, is itself a useful validation: it confirms both that the closed-form expression has been derived and implemented correctly, and that the numerical integration (grid resolution, trapezoidal rule) is fine enough to reproduce the analytic result to well within plotting precision. The qualitative conclusion the extension sets out to demonstrate -- that the particle-in-a-box eigenstates obey, and for low $n$ substantially exceed, the Heisenberg uncertainty principle -- is thus established by two independent routes rather than one.

% =====================================================================
\section{Discussion}
\label{sec:discussion}
% =====================================================================

\subsection{Recurring validation strategy}
Every one of the ten core programs, and both extensions, includes at least one quantitative or structural self-check, and the specific form of that check falls into a small number of recurring categories. Several tasks compare a simulated or numerically evaluated quantity against an independent closed-form theoretical prediction over a full parameter range rather than at a single point: Task~1's mean-squared displacement against $Ns^2$ over all $N$ (Fig.~\ref{fig:task1}), and Task~7's extension comparing a closed-form analytic formula against an independent numerical integration over a full range of $n$ (Fig.~\ref{fig:task7ext}). Others check that a computed quantity correctly recovers a known input parameter: Task~6's linear fit recovering the two input graphite spacings to three decimal places (Fig.~\ref{fig:task6}), and Task~5's computed ground-state energy matching the standard tabulated value of $\SI{13.6057}{eV}$ to seven significant figures. Others check a physically required limiting behaviour: Task~3's Einstein heat capacity saturating at the Dulong--Petit value $3R$ at high temperature (Fig.~\ref{fig:task3b}), Task~9's Compton-scattering quantities correctly vanishing (or reaching a required limiting angle) at $\theta=0$, and Task~4's stopping-voltage lines correctly terminating at the threshold frequency $f_0=\phi/h$ (Fig.~\ref{fig:task4}). Task~8 uses an exact algebraic identity between two competing formulas as its check, isolating the physically meaningful difference (interference cross-term) between them. This variety of validation strategy -- full-curve agreement, parameter recovery, limiting-case behaviour, and algebraic identity -- reflects the different mathematical character of the ten models, and is, in aggregate, a stronger form of validation than any single strategy applied uniformly would have been.

\subsection{From classical statistics to full quantum mechanics}
Read in sequence, the ten tasks trace the historical development of quantum theory reasonably closely. Tasks~1--2 are purely classical (a statistical random walk and a Newtonian hard-sphere collision simulation) and establish the diffusive/statistical baseline against which quantum effects are later a departure. Task~3 introduces energy quantisation for the first time, in its two earliest historical contexts (radiation and lattice vibrations), while retaining a largely classical statistical-mechanics framework (the equilibrium Bose-like occupation number, and the classical equipartition limit recovered at high $T$). Tasks~4--6 are the core results of the ``old quantum theory'' (1900--1924): light quanta (Task~4), the semiclassical Bohr atom (Task~5), and wave-particle duality for matter (Task~6) -- each combining one quantised or discretised assumption with otherwise entirely classical mechanics. Task~7 is the first program in the collection to solve, and evaluate exactly, a genuine time-independent Schr\"odinger-equation eigenvalue problem, and its extension is the first to test a purely quantum-mechanical structural constraint (the uncertainty principle) with no classical analogue. Task~8 addresses quantum measurement itself, contrasting a hidden-variable classical prediction with the Born-rule quantum prediction. Task~9 returns to a scattering problem, but treated fully relativistically, appropriate to its higher photon energy. Task~10 completes the sequence with the exact three-dimensional hydrogenic wavefunctions, the fully quantum-mechanical generalisation of the semiclassical Bohr picture of Task~5.

\subsection{Common numerical idealisations}
A small number of simplifying choices recur across multiple tasks and are worth stating together rather than only within each task's individual limitations paragraph. First, several tasks (4, 5, 7, 10) treat the electron non-relativistically, appropriate to the low energies involved (a few eV to a few hundred eV in the bound-state problems) but explicitly not used in Task~9, where the photon energy is comparable to $m_ec^2/5$ and a full relativistic treatment is essential; Task~6 similarly uses a non-relativistic electron energy--momentum relation, which is a reasonable approximation only because the accelerating voltages used ($1$--$\SI{5}{kV}$) are small compared with the electron rest energy ($\SI{511}{keV}$). Second, several tasks idealise material properties as single sharp numbers -- a single work function per metal (Task~4), a single characteristic Einstein/Debye temperature per solid (Task~3), a single pair of lattice spacings for graphite (Task~6) -- where a real material has some intrinsic spread or additional structure. Third, two tasks (2, 4-extension) use an explicitly acknowledged phenomenological rather than first-principles element -- Task~2's deterministic collision-timing interval in place of a true Poissonian free-path distribution, and the Task~4 extension's uniformly sampled photoelectron speed in place of a physically derived energy distribution -- in both cases because the task specifically calls for an illustrative/animated model rather than a first-principles microscopic simulation.

% =====================================================================
\section{Conclusions}
\label{sec:conclusions}
% =====================================================================

Ten computational physics programs, and two further interactive/verification extensions, have been documented, their underlying physics and mathematics derived from first principles, and their numerical output validated wherever an independent analytical check was available. In every case, the programs' output was found to be consistent with the relevant theory to the precision expected of a double-precision floating-point implementation: the random walk of Task~1 reproduces the diffusive scaling law $\langle r^2\rangle=Ns^2$ to within sampling error; the Einstein heat capacity of Task~3 correctly saturates at the Dulong--Petit limit; the Bohr model of Task~5 reproduces the hydrogen ground-state energy to seven significant figures; the electron-diffraction fit of Task~6 recovers both input graphite spacings to three decimal places; the closed-form uncertainty product of the Task~7 extension agrees with an independent numerical integration of the wavefunction at every quantum number tested, including the ground state; the Compton-scattering kinematics of Task~9 satisfy the required forward-scattering limit and transverse-momentum-conservation identities; and the hydrogenic-orbital angular structure of Task~10 was confirmed by direct evaluation of the predicted equatorial-to-polar density ratio for the $3d_{z^2}$ state. This consistent pattern -- every model checked quantitatively against an independent calculation, rather than accepted merely because it produces a plausible-looking plot -- is the central methodological aim of the collection: to demonstrate not only that each program runs and produces a figure, but that the figure it produces is the one the underlying physics actually predicts.

% =====================================================================
\begin{thebibliography}{99}

\bibitem{codata2018} E.~Tiesinga, P.~J.~Mohr, D.~B.~Newell and B.~N.~Taylor, ``CODATA Recommended Values of the Fundamental Physical Constants: 2018,'' \emph{Rev. Mod. Phys.} \textbf{93}, 025010 (2021); values as tabulated by NIST, \url{https://physics.nist.gov/cuu/Constants/}.

\bibitem{numpy} C.~R.~Harris \emph{et al.}, ``Array programming with NumPy,'' \emph{Nature} \textbf{585}, 357--362 (2020).

\bibitem{matplotlib} J.~D.~Hunter, ``Matplotlib: A 2D Graphics Environment,'' \emph{Computing in Science \& Engineering} \textbf{9}, 90--95 (2007).

\bibitem{einstein1905} A.~Einstein, ``\"Uber die von der molekularkinetischen Theorie der W\"arme geforderte Bewegung von in ruhenden Fl\"ussigkeiten suspendierten Teilchen,'' \emph{Ann. Phys.} \textbf{322}, 549--560 (1905).

\bibitem{einstein1905photoelectric} A.~Einstein, ``\"Uber einen die Erzeugung und Verwandlung des Lichtes betreffenden heuristischen Gesichtspunkt,'' \emph{Ann. Phys.} \textbf{322}, 132--148 (1905).

\bibitem{hyperphysics} R.~Nave, ``Work Functions for Photoelectric Effect,'' HyperPhysics, Georgia State University, citing D.~Halliday and R.~Resnick and the \emph{CRC Handbook of Chemistry and Physics}, \url{https://hyperphysics.phy-astr.gsu.edu/hbase/Tables/photoelec.html}.

\bibitem{bohr1913} N.~Bohr, ``On the Constitution of Atoms and Molecules, Part~I,'' \emph{Philos. Mag.} \textbf{26}, 1--25 (1913).

\bibitem{nist_hydrogen} NIST Atomic Spectra Database / \emph{Handbook of Basic Atomic Spectroscopic Data}, Hydrogen (H~I), \url{https://physics.nist.gov/PhysRefData/Handbook/Tables/hydrogentable1.htm}.

\bibitem{kittel} C.~Kittel, \emph{Introduction to Solid State Physics}, 8th ed. (Wiley, 2004).

\bibitem{graphite_spacing} International Tables for Crystallography / standard graphite structure data, interplanar spacings $d_{002}\approx\SI{0.335}{nm}$ (interlayer) and in-plane reflections consistent with the $\SI{0.123}{nm}$ and $\SI{0.213}{nm}$ spacings used in the standard graphite electron-diffraction tube demonstration.

\bibitem{griffiths} D.~J.~Griffiths and D.~F.~Schroeter, \emph{Introduction to Quantum Mechanics}, 3rd ed. (Cambridge University Press, 2018).

\bibitem{nielsen_chuang} M.~A.~Nielsen and I.~L.~Chuang, \emph{Quantum Computation and Quantum Information}, 10th anniversary ed. (Cambridge University Press, 2010).

\bibitem{aspect1982} A.~Aspect, J.~Dalibard and G.~Roger, ``Experimental Test of Bell's Inequalities Using Time-Varying Analyzers,'' \emph{Phys. Rev. Lett.} \textbf{49}, 1804--1807 (1982).

\bibitem{compton1923} A.~H.~Compton, ``A Quantum Theory of the Scattering of X-rays by Light Elements,'' \emph{Phys. Rev.} \textbf{21}, 483--502 (1923).

\end{thebibliography}

% =====================================================================
\appendix
\section{Source Code}
\label{sec:appendix}
% =====================================================================

This appendix reproduces, in full and unmodified, the source code of every program discussed in this paper. All programs were executed exactly as listed to generate the figures and printed numerical values quoted in the main text.

\subsection{Task 1: Random Walk (\texttt{task1.py})}
\lstinputlisting[language=Python]{code/task1.py}

\subsection{Task 2: Brownian Motion (\texttt{task2.py})}
\lstinputlisting[language=Python]{code/task2.py}

\subsection{Task 3: Planck Spectrum and Einstein Solid (\texttt{task3.py})}
\lstinputlisting[language=Python]{code/task3.py}

\subsection{Task 4: Photoelectric Effect (\texttt{task4.py})}
\lstinputlisting[language=Python]{code/task4.py}

\subsection{Task 4 Extension: Interactive Photoelectric Simulator (\texttt{task4\_extension.py})}
\lstinputlisting[language=Python]{code/task4_extension.py}

\subsection{Task 5: Bohr Model and Hydrogen Spectrum (\texttt{task5.py})}
\lstinputlisting[language=Python]{code/task5.py}

\subsection{Task 6: Electron Diffraction (\texttt{task6.py})}
\lstinputlisting[language=Python]{code/task6.py}

\subsection{Task 7: Particle in a Box (\texttt{task7.py})}
\lstinputlisting[language=Python]{code/task7.py}

\subsection{Task 7 Extension: Uncertainty Principle (\texttt{task7\_extension.py})}
\lstinputlisting[language=Python]{code/task7_extension.py}

\subsection{Task 8: Photon Polarisation / Quantum Cryptography (\texttt{task8.py})}
\lstinputlisting[language=Python]{code/task8.py}

\subsection{Task 9: Compton Scattering (\texttt{task9.py})}
\lstinputlisting[language=Python]{code/task9.py}

\subsection{Task 10: Hydrogenic Orbitals (\texttt{task10.py})}
\lstinputlisting[language=Python]{code/task10.py}

\end{document}
